\clearpage
\chap{A bounded repair budget can recover existence}\label{sec:bounded-realization}
The preceding example does not merely lack a convergent iteration. It has no bounded family of increasingly accurate protected corrections. This suggests a converse question. If approximate corrections stay in one bounded region, does an exact answer exist?

For bounded linear equations in a Hilbert space, the answer is yes. This is an existence statement, not a claim that the given sequence converges in norm. A stronger conclusion becomes available when the candidates are the least norm solutions of nested finite constraint systems.

\begin{scopebox}
This chapter and the next three retain the WR01--WR12 identifiers. Every proof is printed. WR01 and WR02 now have complete formal counterparts in the accepted version 1.8.0 combined Lean project. The printed proof of WR01 uses the weak-subsequence lemma below, which remains a written background argument. Acceptance of the formal counterpart does not by itself verify that printed proof route. The later regularized optimizer and parameter results WR06--WR07 are also included in the accepted combined project. Fixed-parameter estimates within WR09 are formalized, but its limit statement is not. Version 1.10.0 adds complete formal counterparts for WR03--WR05: the attained budget profile, its sharp dual certificate, and its robustness on a fixed correction space. Each result carries its own evidence label. Exact finite models have separate executable checks. The Hilbert and regularization background is classical \cite{clason-inverse}. No novelty is claimed for weak compactness or quadratic regularization.
\end{scopebox}

\subsection{The correction space is the protected kernel}
Let $E,Y,Z$ be real Hilbert spaces, let $d:E\to Y$ and $p:E\to Z$ be bounded linear maps, and fix a raw candidate $f\in E$. Write
\[
 H=\ker p,\qquad T=d|_H:H\to Y,\qquad r=df.
\]
The space $H$ is a closed Hilbert subspace with the norm inherited from $E$. An allowed correction is $a\in H$, and its repaired candidate is $f-a$. Thus
\[
 p(f-a)=pf,\qquad d(f-a)=r-Ta.
\]
The protected exact problem is precisely $Ta=r$. Restricting the domain before optimizing is important. A penalty on $\norm{pa}$ would only discourage a change in the protected data, not prohibit it.

The adjoint below always means $T^*:Y\to H$, taken with the inherited Hilbert metric on $H$. In ambient coordinates this is the projection of $d^*$ onto $H$. A coordinate basis of $\ker p$ need not be orthonormal. Its Gram matrix must then be included in any calculation of correction energy.

\subsection{The compactness needed for the existence argument}
\begin{lemma}[Bounded sequences have weakly convergent subsequences]\label{lem:weak-subsequence}
Every bounded sequence $(a_n)$ in a real Hilbert space has a subsequence converging weakly to some $a$. Moreover,
\[
 \norm a\le\liminf_k\norm{a_{n_k}}.
\]
\end{lemma}
\begin{proof}
Choose $M\ge0$ with $\norm{a_n}\le M$. If $M=0$, every term is zero and the conclusion is immediate. Otherwise let $H_0$ be the closed linear span of the countably many vectors $a_n$. Finite rational linear combinations of those vectors form a countable dense set, so $H_0$ is separable. Apply Gram--Schmidt to a countable dense sequence, omitting zero remainders. This gives either a finite orthonormal basis or a countable orthonormal sequence $(e_j)$ whose closed span is $H_0$. The ambient Hilbert space itself need not be separable.

Each scalar sequence $\langle a_n,e_j\rangle$ is bounded by $M$. For the first coordinate choose a convergent subsequence, for the second choose a subsequence of that subsequence, and continue. Taking the $k$th term from the $k$th nested subsequence, with increasing original indices, yields $a_{n_k}$ for which every fixed coordinate converges. Call the coordinate limits $c_j$. In the finite-dimensional case only finitely many successive selections are required.

For each finite $J$, Bessel's inequality and passage to the limit give
\[
 \sum_{j=1}^Jc_j^2
 =\lim_k\sum_{j=1}^J\langle a_{n_k},e_j\rangle^2\le M^2.
\]
The nondecreasing partial sums are therefore bounded. Orthogonality makes the partial vector sums $\sum_{j\le J}c_je_j$ Cauchy, so completeness defines $a\in H_0$ with $\norm a\le M$.

To verify weak convergence against $h\in H_0$, choose a finite basis combination $h_J$ close to $h$. Then
\[
 |\langle a_{n_k}-a,h\rangle|
 \le |\langle a_{n_k}-a,h_J\rangle|
      +2M\norm{h-h_J}.
\]
For a prescribed tolerance, choose $h_J$ first to control the second term and then $k$ to control the first by finitely many coordinate limits. For a general $h$ in the ambient space, write $h=P_{H_0}h+h_\perp$. The orthogonal component pairs to zero with every $a_{n_k}$ and with $a$, reducing to the preceding case. This proves weak convergence against every test vector.

If $a\ne0$, testing against $a/\norm a$ gives
\[
 \norm a=\lim_k\langle a_{n_k},a/\norm a\rangle
 \le\liminf_k\norm{a_{n_k}}.
\]
The zero limit case follows from nonnegativity. Thus the lower-semicontinuity conclusion concerns the actual norms on the selected subsequence, not only the original bound $M$.
\end{proof}

\begin{theorem}[Bounded approximate feasibility]\label{WR01}
For a bounded linear map $T:H\to Y$ between real Hilbert spaces and $r\in Y$, the following are equivalent:
\[
 \begin{split}
 &r\in\ran T,\\
 &\text{there is a bounded sequence }(a_n)\subset H
       \text{ with }\norm{Ta_n-r}\longrightarrow0.
 \end{split}
\]
If all $\norm{a_n}\le M$, an exact solution exists with norm at most $M$.
\end{theorem}
\begin{proof}
If $Ta=r$, the constant sequence $a_n=a$ is bounded and has zero residual. For the converse, let $\norm{a_n}\le M$ and $Ta_n\to r$ in norm. By the preceding lemma there is a subsequence $a_{n_k}$ converging weakly to $a\in H$ with $\norm a\le M$.

For every fixed $y\in Y$, the adjoint identity and weak convergence give
\[
 \langle Ta,y\rangle
 =\langle a,T^*y\rangle
 =\lim_k\langle a_{n_k},T^*y\rangle
 =\lim_k\langle Ta_{n_k},y\rangle.
\]
The last limit equals $\langle r,y\rangle$ because norm convergence implies convergence of each continuous pairing. Thus $\langle Ta-r,y\rangle=0$ for every $y$. Taking $y=Ta-r$ gives $\norm{Ta-r}^2=0$, hence $Ta=r$.

This proof does not claim that $Ta_{n_k}$ converges in norm merely because $a_{n_k}$ converges weakly. Its norm convergence to $r$ was a separate hypothesis. The boundedness of the candidate corrections is used to obtain a weakly convergent subsequence, and their accurate residuals are used to identify its image exactly. The bound $\norm a\le M$ comes from weak lower semicontinuity, so the exact answer stays inside the same closed budget ball.
\end{proof}
\written{WR01}

No closed range assumption appears. The theorem concerns a single datum and a bounded family of its corrections. Closed range is the stronger, uniform statement that every attainable datum admits a correction controlled by one residual gain.

An immediate contrapositive is useful. If $r\in\overline{\ran T}\setminus\ran T$, then every sequence with $Ta_n\to r$ satisfies $\norm{a_n}\to\infty$. Otherwise infinitely many terms would lie in one bounded ball, contradicting the theorem along that subsequence. This recovers the escape mechanism of Section~\ref{sec:escape} without using the diagonal formula.

\subsection{Existence does not make an arbitrary sequence converge strongly}
Take $H=Y=\ell^2$ and $Te_n=e_n/n$. The vectors $a_n=e_n$ have norm one and satisfy $Ta_n\to0$. An exact answer exists, namely zero. Yet $\norm{e_n-e_m}=\sqrt2$ for $n\ne m$, so no subsequence converges in norm.

The conclusion of WR01 is the existence of a weak limit along a subsequence that solves the equation. It is not summable movement, a strong convergence rate, or a selection rule for the original candidates. The minimum correction for the zero datum is zero, which this particular sequence never approaches in norm.

\subsection{Minimum energy finite repairs do converge strongly}
\fixedplate{31_nested_energy}{Minimum energy finite corrections behave differently for two data under the same diagonal operator. Bounded partial energy produces a strong Hilbert realization. Escaping partial energy accompanies a coherent product-space point outside the Hilbert space. The attainable energy lies below the explicit bound two.}{fig:nested-energy}

Suppose now that $A_1\supseteq A_2\supseteq\cdots$ are nonempty closed affine subsets of a real Hilbert correction space $H$. A typical $A_n$ consists of corrections satisfying the first $n$ bounded linear observation equations. The sets are affine because the equations are linear, and closed because the observations are continuous.

Let $a_n$ be the unique point of $A_n$ nearest to zero, and put $m_n=\norm{a_n}$. The Hilbert projection theorem supplies these points. It can also be proved directly: a sequence approaching the minimum norm is Cauchy by the parallelogram identity, its limit belongs to the closed convex set, and the same identity gives uniqueness.

\begin{theorem}[A uniform finite budget yields a strong realization]\label{WR02}
For the nested affine sets above, $m_n$ is nondecreasing, and for $m\ge n$,
\begin{equation}\label{eq:nested-energy}
 \norm{a_m-a_n}^2=m_m^2-m_n^2.
\end{equation}
Their intersection is nonempty if and only if $\sup_n m_n<\infty$. In that case $a_n$ converges in norm to the unique least norm point $a_*\in\bigcap_n A_n$, with
\[
 \norm{a_*-a_n}^2=\norm{a_*}^2-\norm{a_n}^2.
\]
If the intersection is empty, then $m_n\to\infty$.
\end{theorem}
\begin{proof}
The nearest-point theorem for nonempty closed convex subsets of a Hilbert space, proved in Appendix~\ref{app:proof-tools}, supplies $a_n$. Write $A_n=a_n+L_n$, where $L_n$ is its closed direction subspace. For $\ell\in L_n$ and every real $t$, the vector $a_n+t\ell$ belongs to $A_n$. Minimality implies
\[
 0\le\norm{a_n+t\ell}^2-\norm{a_n}^2
 =2t\langle a_n,\ell\rangle+t^2\norm\ell^2.
\]
Allowing both signs of arbitrarily small $t$ forces $\langle a_n,\ell\rangle=0$.

If $m\ge n$, nesting gives $a_m\in A_n$, so $a_m-a_n\in L_n$. Expanding its norm therefore gives
\[
 \norm{a_m}^2
 =\norm{a_n+(a_m-a_n)}^2
 =\norm{a_n}^2+\norm{a_m-a_n}^2.
\]
This is the exact energy identity. It also makes $m_n^2$ nondecreasing.

Suppose these energies are bounded, and let $L=\lim_n m_n^2$. Given $\epsilon>0$, choose $N$ with $L-m_N^2<\epsilon^2$. For $m\ge n\ge N$,
\[
 \norm{a_m-a_n}^2=m_m^2-m_n^2\le L-m_N^2<\epsilon^2.
\]
Thus $(a_n)$ is Cauchy and has a norm limit $a_*$. For each fixed $n$, all terms from that index onward belong to the closed set $A_n$, so $a_*\in A_n$. This places $a_*$ in their entire intersection.

If $a$ lies in the intersection, then $\norm a\ge\norm{a_n}$ for every $n$. Norm continuity gives $\norm a\ge\norm{a_*}$. Uniqueness follows from the nearest-point theorem applied to the intersection, which is closed and convex and now known to be nonempty. Alternatively, the midpoint of two different minimizers would have strictly smaller norm by the parallelogram identity. Taking $m\to\infty$ in the finite energy identity yields
\[
 \norm{a_*-a_n}^2=\norm{a_*}^2-\norm{a_n}^2.
\]
Conversely, any element of the intersection bounds every minimum norm from above. Therefore an empty intersection forces the nondecreasing sequence $m_n$ to be unbounded. A nondecreasing unbounded nonnegative sequence tends to infinity, giving the last conclusion.
\end{proof}
\written{WR02}

The finite answers need not extend one another coordinate by coordinate. Their shared feature is minimum energy in one common ambient Hilbert space. This is the extra geometry that forces their differences to be orthogonal to the earlier minimizers.

For the attainable diagonal datum $r_n=1/n^2$ under $Tu=(u_n/n)$, with $n\ge1$, the first $N$ equations force correction coordinate $n$ to equal $1/n$ for $n\le N$. Their minimum energy completion is zero beyond $N$, and
\[
 m_N^2=\sum_{n=1}^{N}\frac1{n^2}\le2.
\]
These corrections converge strongly to $(1/n)$. For the escaping datum $r_n=1/n$, the corresponding energy is $m_N^2=N$. Every finite equation is again solved, but the budget test now rules out an ambient solution.


\begin{insight}
Finite consistency is not enough by itself. For nested closed affine problems in one Hilbert space, a uniform bound on the minimum correction norms is exactly what upgrades finite feasibility to ambient feasibility. This is a datum-specific test, not a uniform inverse bound for the operator.
\end{insight}
